\section{Normal forms}
\label{sec:prop-syntax-semantics-04}

\begin{definition}[Propositional literal, clause, and formula interpretation]
A \emph{propositional literal} is a propositional variable or the negation of
a propositional variable.  A \emph{clause} is a finite set of propositional
literals.  Thus clause equality is set equality: order and repeated
occurrences are not part of a clause.

To map this set object back to the formula language, put
\[
 \operatorname{code}(p_n)=2n,
 \qquad
 \operatorname{code}(\neg p_n)=2n+1.
\]
For a clause \(C\), let \(k=\abs{C}\); when \(k>0\), list its members as
\(L_1,\ldots,L_k\) in strictly increasing code order and define, by downward
recursion,
\[
 D_k\defeq L_k,
 \qquad
 D_i\defeq(L_i\lor D_{i+1})\quad(1\leq i<k).
\]
Then define
\[
 \operatorname{form}_{\CPL}(C)\defeq
 \begin{cases}
  \bottom,&C=\varnothing,\\
  D_1,&C\ne\varnothing.
 \end{cases}
\]
This is an exact metalinguistic function from clauses to
\(\Fm_{\CPL}\), not a new formula constructor.  A valuation \(v\)
\emph{satisfies} \(C\) iff
\(v\entails\operatorname{form}_{\CPL}(C)\).  In particular, the empty
clause is false under every valuation.
\end{definition}

\begin{definition}[Negation, conjunctive, and disjunctive normal form]
A formula is in \emph{negation normal form} iff it is built from literals and
\(\bottom,\topv\) using only \(\land\) and \(\lor\).  It is in
\emph{conjunctive normal form} iff it is \(\topv\), the empty conjunction,
or there are clauses \(C_1,\ldots,C_m\), with \(m>0\), such that it is
\(N_1\), where the formula sequence \(N_m,\ldots,N_1\) is defined by
\[
 N_m\defeq\operatorname{form}_{\CPL}(C_m),
 \qquad
 N_i\defeq
 \bigl(\operatorname{form}_{\CPL}(C_i)\land N_{i+1}\bigr)
 \quad(1\leq i<m).
\]
Thus a conjunctive-normal-form formula is a right-associated finite
conjunction of the formula interpretations of clauses.  In every finite
conjunction and finite disjunction used in this chapter, the empty conjunction
is \(\topv\) and the empty disjunction is \(\bottom\).  A formula is in
\emph{disjunctive normal form} iff it is a finite disjunction of finite
conjunctions of literals under these conventions.
\end{definition}

\begin{lemma}[Satisfaction of clauses and of clause conjunctions]
\label{lem:prop-clause-satisfaction}
Let \(v\) be a valuation.
\begin{enumerate}
  \item For every clause \(C\):
        \(v\entails\operatorname{form}_{\CPL}(C)\) iff
        \(v\entails L\) for at least one literal \(L\in C\).
  \item For all clauses \(C_1,\ldots,C_m\) with \(m>0\), and the formulas
        \(N_m,\ldots,N_1\) built from them as in the definition of
        conjunctive normal form: \(v\entails N_1\) iff
        \(v\entails\operatorname{form}_{\CPL}(C_j)\) for every \(j\) with
        \(1\leq j\leq m\).
\end{enumerate}
\end{lemma}
\begin{proof}
For 1, if \(C=\varnothing\), then
\(\operatorname{form}_{\CPL}(C)=\bottom\), which no valuation satisfies,
and \(C\) has no member, so both sides fail.  Otherwise let
\(k=\abs{C}>0\), list \(C\) in strictly increasing code order as
\(L_1,\ldots,L_k\), and let \(D_k,\ldots,D_1\) be the formulas of the
clause-interpretation recursion.  By downward induction along that
recursion, \(v\entails D_i\) iff \(v\entails L_j\) for some \(j\) with
\(i\leq j\leq k\): the base case \(D_k=L_k\) is immediate, and the case
\(D_i=(L_i\lor D_{i+1})\) follows from the recursive truth clause for
\(\lor\) together with the induction hypothesis.  Taking \(i=1\) and
\(\operatorname{form}_{\CPL}(C)=D_1\) proves 1.

For 2, by downward induction along the recursion defining
\(N_m,\ldots,N_1\), \(v\entails N_i\) iff
\(v\entails\operatorname{form}_{\CPL}(C_j)\) for every \(j\) with
\(i\leq j\leq m\): the base case \(N_m=\operatorname{form}_{\CPL}(C_m)\)
is immediate, and the case
\(N_i=(\operatorname{form}_{\CPL}(C_i)\land N_{i+1})\) follows from the
recursive truth clause for \(\land\) together with the induction
hypothesis.  Taking \(i=1\) proves 2.
\end{proof}

\begin{theorem}[Negation normal form]
Every propositional formula is logically equivalent to a formula in negation
normal form.
\end{theorem}
\begin{proof}
First eliminate every occurrence of \(\to\) using
\(\varphi\to\psi\equiv\neg\varphi\lor\psi\).  Proceed by structural
induction through the resulting formula.  Double negations are removed, and a
negation immediately above a binary connective is moved inward by
\[
 \neg(\varphi\land\psi)\equiv\neg\varphi\lor\neg\psi,
 \qquad
 \neg(\varphi\lor\psi)\equiv\neg\varphi\land\neg\psi.
\]
The base cases are variables, their negations, \(\bottom\), and \(\topv\).
Each displayed replacement decreases the number of connectives lying below a
negation, so the process terminates.  Every replacement is checked by the
recursive truth clauses and therefore preserves logical equivalence.
\end{proof}

\begin{theorem}[Conjunctive and disjunctive normal forms]
\label{thm:prop-normal-forms}
Every propositional formula is logically equivalent to a formula in
disjunctive normal form and to a formula in conjunctive normal form.
\end{theorem}
\begin{proof}
Let \(p_1,\ldots,p_n\) be the variables occurring in \(\varphi\), listed in
increasing index order.  In this
proof only, a \emph{valuation on} a finite variable set is a function from
that set to \(\B\); the recursive truth clauses apply verbatim to formulas
whose variables lie in the set, and \cref{lem:prop-coincidence} identifies
satisfaction under such a function with satisfaction under any total
valuation extending it.  For each
valuation \(v\) on these variables define
\[
  \delta_v\defeq\bigwedge_{i=1}^{n}\ell_i^v,
  \qquad
  \ell_i^v\defeq
  \begin{cases}
    p_i,&v(p_i)=\True,\\
    \neg p_i,&v(p_i)=\False.
  \end{cases}
\]
A valuation on the empty variable set is the empty function, and when
\(n=0\) the displayed conjunction is \(\topv\) by the preceding convention.
A valuation \(w\) satisfies \(\delta_v\) iff \(w(p_i)=v(p_i)\) for every
\(i\).  Let \(V_{\varphi}\) be the finite set of valuations on
\(p_1,\ldots,p_n\) that satisfy \(\varphi\).  Order these valuations
lexicographically by
\((v(p_1),\ldots,v(p_n))\), with \(\False<\True\), enumerate the resulting
order as \(v_1,\ldots,v_m\), and put
\[
  D_{\varphi}\defeq\bigvee_{j=1}^{m}\delta_{v_j}.
\]
If \(V_{\varphi}=\varnothing\), then \(m=0\) and this disjunction is
\(\bottom\).  By the
coincidence lemma, an arbitrary valuation satisfies \(D_{\varphi}\) iff
it satisfies \(\varphi\).  Hence \(D_{\varphi}\) is the required
disjunctive normal form.

For the conjunctive normal form, construct dually.  For each valuation
\(v\) on \(p_1,\ldots,p_n\) put
\[
  \bar\ell^{\,v}_i\defeq
  \begin{cases}
    \neg p_i,&v(p_i)=\True,\\
    p_i,&v(p_i)=\False,
  \end{cases}
  \qquad
  K_v\defeq\set{\bar\ell^{\,v}_1,\ldots,\bar\ell^{\,v}_n}.
\]
The coding in the clause-interpretation definition gives the variable of
index \(i\) the code \(2i\) and its negation the code \(2i+1\), both
strictly increasing in \(i\).  The enumeration \(p_1,\ldots,p_n\) is in
increasing index order---the indices need not form an initial
segment---so any literal on \(p_j\) has smaller code than any literal on
\(p_k\) for \(j<k\): if their indices are \(i_j<i_k\), then
\(2i_j+1<2i_k\).  The listing of \(K_v\) in strictly increasing code order
is therefore exactly \(\bar\ell^{\,v}_1,\ldots,\bar\ell^{\,v}_n\), and
\(\gamma_v\defeq\operatorname{form}_{\CPL}(K_v)\) requires no
re-ordering.  A valuation \(w\) falsifies \(\bar\ell^{\,v}_i\) iff
\(w(p_i)=v(p_i)\), so \cref{lem:prop-clause-satisfaction}(1) gives:
\(w\) falsifies \(\gamma_v\) iff \(w(p_i)=v(p_i)\) for every
\(i\leq n\).

Let \(W_\varphi\) be the set of valuations on \(p_1,\ldots,p_n\) that
falsify \(\varphi\), enumerated in the lexicographic order fixed above as
\(v_1,\ldots,v_m\).  Put \(C_\varphi\defeq\topv\) if \(m=0\); otherwise
put \(C_\varphi\defeq N_1\), where \(N_m,\ldots,N_1\) are built from the
clauses \(C_j\defeq K_{v_j}\) by the recursion in the definition of
conjunctive normal form.  In every case \(C_\varphi\) satisfies that
definition verbatim, the degenerate cases included: for valid
\(\varphi\), \(m=0\) and \(C_\varphi=\topv\); for variable-free true
\(\varphi\), \(n=m=0\) and again \(C_\varphi=\topv\); for variable-free
false \(\varphi\), \(n=0\) and \(m=1\) with \(v_1\) the empty valuation,
so \(K_{v_1}=\varnothing\) and
\(C_\varphi=\operatorname{form}_{\CPL}(\varnothing)=\bottom\), the
conjunctive normal form whose single clause is empty.

Let \(w\) be an arbitrary valuation and \(w_0\) its restriction to
\(\set{p_1,\ldots,p_n}\).  Suppose \(m>0\).  By
\cref{lem:prop-clause-satisfaction}(2) and (1), \(w\) satisfies
\(C_\varphi\) iff for every \(j\leq m\) it satisfies some literal of
\(K_{v_j}\), iff \(w_0\ne v_j\) for every \(j\leq m\), iff
\(w_0\notin W_\varphi\), iff \(w_0\) satisfies \(\varphi\), iff, by the
coincidence lemma, \(w\) satisfies \(\varphi\).  If \(m=0\), every
valuation on the variable set satisfies \(\varphi\), so by the
coincidence lemma every valuation satisfies both \(\varphi\) and
\(C_\varphi=\topv\).  Hence \(C_\varphi\) is the required conjunctive
normal form, obtained from the definitional recursions,
\cref{lem:prop-clause-satisfaction}, and the coincidence lemma alone.
\end{proof}

For the satisfiability-preserving transformation with fresh variables,
see Section~\ref{sec:p02-tseitin-transformation} and
Theorem~\ref{thm:p02-tseitin}.

An \(n\)-ary \emph{truth-functional connective} is a connective symbol
interpreted by a fixed function \(\B^n\to\B\).  For a set \(C\) of
truth-functional connectives, the \emph{formulas over \(C\)} are generated
by the propositional variables together with one formation clause per member
of \(C\), and the recursive truth clause of each member applies its
interpreting function to the values of the immediate subformulas.  For a
formula \(\varphi\) whose variables are among \(p_1,\ldots,p_n\), the
\emph{truth function of \(\varphi\) in those variables} is the function
\(f_\varphi:\B^n\to\B\) with
\(f_\varphi(a_1,\ldots,a_n)\defeq\widehat v(\varphi)\) for any valuation
\(v\) with \(v(p_i)=a_i\) for each \(i\); this is well defined by
\cref{lem:prop-coincidence}.

\begin{definition}[Functional completeness]
A set \(C\) of truth-functional connectives is \emph{functionally complete}
iff, for every \(n\in\N\) and every function \(f:\B^n\to\B\), some formula
built from variables \(p_1,\ldots,p_n\) and connectives in \(C\) has truth
function \(f\) in those variables.
\end{definition}

\begin{theorem}[A functionally complete basis]
The set \(\set{\neg,\land,\bottom}\) is functionally complete.
\end{theorem}
\begin{proof}
Given \(f:\B^n\to\B\), form the disjunctive normal form whose conjunctions of
literals correspond exactly to the inputs \(a\in\B^n\) for which
\(f(a)=\True\), as in \cref{thm:prop-normal-forms}.  Replace each occurrence
of \(\alpha\lor\beta\) by
\(\neg(\neg\alpha\land\neg\beta)\).  The empty disjunction is \(\bottom\).
The resulting formula uses only \(\neg,\land,\bottom\) and computes \(f\).
\end{proof}
